\documentclass{beamer}
\usetheme{Madrid}
\usepackage{amsmath, amssymb, listings, xcolor}

\title{Understanding Pipelining Hazards}
\subtitle{BLG 483E -- Artificial Intelligence}
\author{salihsefer36}
\date{\today}

\begin{document}

% --- Slide 1: Title ---
\begin{frame}
  \titlepage
\end{frame}

% --- Slide 2: Introduction to CPU Pipelining ---
\begin{frame}{Introduction to CPU Pipelining}
  \textit{Think of a restaurant: one waiter takes orders, one chef cooks,
  another plates, another serves. If one person did everything per customer,
  the restaurant would serve one table at a time. A CPU pipeline works
  exactly the same way.}

  \medskip
  \begin{itemize}
    \item A pipeline splits every instruction into 5 stages that run
          \textbf{simultaneously} for different instructions:
          \textbf{IF} (fetch), \textbf{ID} (decode), \textbf{EX} (execute),
          \textbf{MEM} (memory), \textbf{WB} (write result).
    \item When the assembly line runs smoothly, the CPU completes roughly
          one instruction per clock cycle (instead of one every 5 cycles).
    \item A \textbf{hazard} is anything that breaks the smooth flow ---
          a conflict that forces a stage to wait or throw away work.
    \item Hazards come in three flavours: \textbf{data}, \textbf{structural},
          and \textbf{control}.
  \end{itemize}
\end{frame}

% --- Slide 3a: Data Hazards — Types ---
\begin{frame}{Data Hazards --- Types}
  \textit{Imagine a relay race where runner B must grab the baton from runner A
  before sprinting. If B takes off before A arrives, the handoff fails.
  Data hazards are exactly that: one instruction needs a value the previous
  one has not finished computing yet.}

  \medskip
  \begin{itemize}
    \item \textbf{RAW} (Read-After-Write, ``true dependence''): the most common
          type. Instruction B tries to read a register before instruction A has
          written its result into it. \textit{Runner B starts before the baton
          arrives.}
    \item \textbf{WAR} (Write-After-Read, ``anti-dependence''): B wants to
          overwrite a register that A is still reading. \textit{Chef B wants to
          clean the pan that chef A is still using.}
    \item \textbf{WAW} (Write-After-Write, ``output dependence''): two
          instructions both write the same register; the second write must
          ``win''. \textit{Two workers submit the same report; only the later
          version should be filed.}
  \end{itemize}
\end{frame}

% --- Slide 3b: Data Hazards — Detection & Mitigation ---
\begin{frame}{Data Hazards --- Detection \& Mitigation}
  \textit{Think of a traffic officer watching an intersection. When two cars
  are about to collide, the officer either holds one back (stall) or waves it
  through a special shortcut lane (forwarding).}

  \medskip
  \begin{itemize}
    \item \textbf{Detection:} A hazard detection unit (``the traffic officer'')
          checks the pipeline registers (ID/EX, EX/MEM, MEM/WB) every cycle
          to spot conflicts before they cause wrong results.
    \item \textbf{Stall (interlock):} Freeze the fetch and decode stages,
          insert a bubble (an empty ``do nothing'' slot) --- like holding a car
          at a red light for one cycle.
    \item \textbf{Forwarding (bypassing):} Route the computed result directly
          from one pipeline stage to the next instruction's input, skipping the
          detour through the register file. Eliminates most RAW stalls completely.
    \item \textbf{Load-use exception:} When instruction A loads from memory and
          B immediately needs that value, forwarding still cannot help --- one
          stall cycle is unavoidable.
  \end{itemize}
\end{frame}

% --- Slide 4: Structural Hazards ---
\begin{frame}{Structural Hazards}
  \textit{Two delivery trucks arriving at a warehouse with a single loading dock
  at the same time: one must wait. Structural hazards are the same idea ---
  two pipeline stages need the same piece of hardware simultaneously.}

  \medskip
  \begin{itemize}
    \item \textbf{Memory conflict:} If the CPU has only one memory unit (no
          separate instruction cache and data cache), the Fetch stage and the
          Memory stage cannot both run in the same cycle --- one must stall.
    \item \textbf{Register write-port conflict:} A register file with a single
          write port cannot accept two write-back results in the same cycle.
          \textit{One printer, two print jobs queued at the same instant.}
    \item \textbf{Multi-cycle unit contention:} Slow operations (floating-point
          divide, square root) tie up a shared unit for many cycles, blocking
          other instructions from using it.
    \item \textbf{Fixes:} Add duplicate hardware (two memory ports), share the
          resource across alternating cycles, or let the compiler schedule
          independent work in between.
  \end{itemize}
\end{frame}

% --- Slide 5: Control Hazards ---
\begin{frame}{Control Hazards}
  \textit{You are driving and your GPS says ``turn left'' --- but you are
  already in the right lane three blocks back. You have ``fetched'' the wrong
  road. A branch hazard is the same: the CPU grabs the wrong next instructions
  before it knows which way a branch (an \texttt{if}/\texttt{loop}) goes.}

  \medskip
  \begin{itemize}
    \item A branch's direction and destination are only known at the EX stage,
          so 2 instructions are already fetched for the wrong path --- they must
          be \textbf{flushed} (thrown away), wasting 2 cycles.
    \item \textbf{Static prediction:} Always guess ``branch not taken'' (keep
          driving straight). Simple, but wrong for loops.
    \item \textbf{Dynamic prediction (2-bit counter):} The CPU remembers past
          behaviour in a small table and guesses based on history --- like a GPS
          that learns your usual route.
    \item \textbf{Branch Target Buffer (BTB):} Caches the destination address
          so the CPU can jump immediately without recomputing where to go.
    \item \textbf{Delayed branch:} The slot right after a branch always executes;
          the compiler fills it with a useful instruction instead of wasting it.
  \end{itemize}
\end{frame}

% --- Slide 6: Hazard Mitigation Summary ---
\begin{frame}{Hazard Mitigation Summary}
  \textit{A factory supervisor has three tools: a stop button (stall), a
  conveyor shortcut (forwarding), and a smart pre-planned schedule (compiler
  reordering). Together they keep the assembly line moving even when hiccups
  occur.}

  \medskip
  \begin{itemize}
    \item \textbf{Stall:} Pause the front of the pipeline and insert an empty
          ``bubble'' instruction --- buys time for the slow stage to finish,
          but costs cycles.
    \item \textbf{Forwarding:} A hardware shortcut that hands a result directly
          from one stage to the next instruction's input, cutting most RAW stalls
          to zero extra cycles.
    \item \textbf{Load-use stall:} The one case forwarding cannot fully fix; a
          single wasted cycle is unavoidable unless the compiler reorders code.
    \item \textbf{Branch delay slot:} The instruction after a branch always
          runs; filling it with real work (not a \texttt{NOP}) recovers a cycle
          for free.
    \item \textbf{Compiler scheduling:} By rearranging independent instructions
          at compile time, the compiler can eliminate many hazards before the
          hardware ever sees them.
  \end{itemize}
\end{frame}

% --- Slide 7: References ---
\begin{frame}{References}
  \begin{thebibliography}{9}
    \bibitem{patterson} Patterson, D.\,A. and Hennessy, J.\,L.
      \textit{Computer Organization and Design: The Hardware/Software Interface},
      5th ed. Morgan Kaufmann, 2014.
    \bibitem{hennessy} Hennessy, J.\,L. and Patterson, D.\,A.
      \textit{Computer Architecture: A Quantitative Approach},
      6th ed. Morgan Kaufmann, 2017.
  \end{thebibliography}
\end{frame}

\end{document}
